\head{第12节末：超奇异点、线性截面族与一般中心}
\source{法稿物理第71页末至73页末；英译物理第47页下部至48页末；第十九版\texttt{v\_dot\_middle.tex}旧D17-M036--045}

\audit{V20-FAM-E01}{恢复字句与分清对象}{法稿可读的内容据实恢复；一般推广仍保留其草稿身份}
本段实际逐页查看法稿71--74页、英译47--49页，并对法稿72页的手写插入及73页末行作高分辨率裁图复核。74页和英译49页仅用于确认第13节边界，并不计入本轮已回校正文。

法稿72页的“cela implique encore que”可以读出，旧D17-M037的“it implies”后缺口由此撤销；“pourvu que”译为条件性的“前提是”。恢复的是文字，不是8.7(c)在早先底本中的缺续或任意$m$的唯一性证明。原作者“如有需要，我可以给出更精确的陈述”仍是原作者的话。

法稿72页排$V^1$，英译47页改为$V'$；正文采用$V^1$，说明其英文别写，保留双撇$V''$以及另外定义的$V^2$。$V^2$是$V^1\setminus V''\to\operatorname{Grass}^m$的分歧子预概形，不是按交维数定义的$V^{(2)}$。环境$X_{\operatorname{Grass}^m}$在底本中确实出现，不能悄改成$X^{(m)}$。关于相对维数小1的句子承接前面的纯维、非空片限制；关于“至少两个普通奇异零点合并”的解释没有在原稿中配上一个精确变形命题，不应假装已经给出。

法稿73页还恢复了$S=\Spec(k)$、$X\times_P C$、小写$f$和$T^1$。英译相应漏$(k)$、损坏纤维积、把$f$排成$F$并把$T^1$改为$T'$。旧D17-M041关于基域完全不可恢复的判断已过时；但法稿随后写全局$T$“在$k(\xi)$上几何不可约”的问题确属法稿，并且并非只把基域改为$k$即可解决。

关于关联概形与爆破的关系，本轮承接本版10.2、10.3已经明确的局部诺特和正则列范围。原稿本段仅说“正则地”并引用上述条目，不另造作者没有写出的证明。对于中心处长度$m+1$的正则列，关联条件正是射影丛中泛核的零点条件；已有爆破论证在相同正则性范围内适用。一般非诺特版本的正则性约定继续列为未决事项。

\audit{V20-FAM-E02}{一般线性束的判别集合不必几何不可约}{原稿存在实质反例；不是英文转写造成的错误}
\oldauditsrc{法稿73页中部：“est g\'eom irr\'eductible sur $k(\xi)$, et de codim 1 dans $Q$”；英译48页同样保留不可约性断言}
法稿把$X$假定为$k$上光滑、固有、几何不可约，$n=\dim X\ge m$；对余维$m+1$的一般中心$C$，考虑余维$m$且包含$C$的线性截面。作者稍后补充$f$无分歧和预先作次数至少2的Veronese变换。即使加入这两项后补条件，其判别集合$T$的几何不可约性仍不成立。

\begin{analysisproof}
取代数闭的特征零域$k$，从恒等闭浸入$\mathbf P^1_k\to\mathbf P^1_k$出发，作二次Veronese变换：
\[
\nu_2\colon X=\mathbf P^1_k\longrightarrow\mathbf P^2_k,
\qquad [u:v]\longmapsto [u^2:uv:v^2].
\]
其像是光滑圆锥曲线$x_0x_2=x_1^2$。因此$X$光滑、固有、几何不可约，$f=\nu_2$为闭浸入，特别无分歧；所用Veronese次数为2。取$m=n=1$以及不在曲线上的中心$C=[0:1:0]$。所有通过$C$的直线组成
\[
Q(C)=\mathbf P^1_{[a:c]},\qquad L_{[a:c]}:a x_0+c x_2=0.
\]
拉回$L_{[a:c]}$的方程为$au^2+cv^2=0$。若$ac\ne0$，它在代数闭域上有两个不同的根，故截面是光滑零维概形。若$a=0$或$c=0$，它是单个长度2的点，局部环同构于$k[t]/(t^2)$。于是
\[
T=V(ac)=\{[1:0],[0:1]\}\subset Q(C).
\]
这两个点互不相同，故$T$非几何不可约。两个例外截面各恰有一个普通二重零点：在局部坐标中二次项的系数非零；特征零保证其二阶型非退化。因此本例$T^1=\varnothing$，且作者关于余维1、普通奇点以及每个坏纤维唯一奇点的其他预期在此例中都不造成问题。反例精确针对不可约性。

还须排除“中心不够一般”的解释。对偶平面中直线$a x_0+b x_1+c x_2=0$的非光滑截面判别式为
\[
\Delta=b^2-4ac.
\]
故所有切线构成光滑的对偶圆锥曲线$Z=V(\Delta)$。一个中心$C=[c_0:c_1:c_2]$对应对偶直线
\[
Q(C): a c_0+b c_1+c c_2=0.
\]
在$[a:b:c]\in Z$处，对偶圆锥曲线的切线系数为$[-4c:2b:-4a]$，其系数满足$c_0c_2=c_1^2$。反过来，这些正是原圆锥曲线上的点所给出的对偶切线。因此，当$C$位于原圆锥曲线的非空开补集时，$Q(C)$与$Z$横截于两个不同的几何点。整个这一非空开参数集都给出可约的$T$，足以否定原稿对“足够一般”的$C$所作的不可约性断言。
\end{analysisproof}

正文将错误断言移入V20-FAM007脚注，保留其完整中文含义、反例和作者后补条件，并明确写出不能推出不可约性。其他拟议结论按原顺序保存，没有因这一个反例就一并判假，也没有声称全部已经证明。原稿两用$n$分别表示维数和Veronese指标；正文仅将后者改名为$d$，原记号记录在V20-FAM009。不能借原稿的重名凭空增加“$\dim X\ge2$”来回避曲线反例。

\audit{V20-FAM-E03}{第12节末一般中心的余维界}{给出独立补证；余维界不推出几何不可约性}
\oldauditsrc{法稿73页最后一段及手写末行；英译48页末段}
法稿在印刷段中明确写$Q(C)\cap Z$，英译的第一处$Q(Z)\cap Z$是转写错。法稿手写末行可读为：“此外，若$Z$几何不可约，$Q(C)\cap Z$也同样如此，如果$Z$‘足够一般’。”其中“itou”即“同样”；最后的$Z$没有改猜为$C$。旧D17-M045留下的谓语缺口因此可以恢复，但恢复后的不可约性断言不成立于它最自然的无附加条件解释。

下面只证明印刷段的余维界。此证明是本轮补证，不是原稿漏写证明的复原。

\begin{analysisproof}
设$P=\mathbf P^N_k$，$1\le m\le N-1$，记
\[
G=\operatorname{Grass}^m(P),\qquad
B=\operatorname{Grass}^{m+1}(P).
\]
$G$参数化余维$m$的线性子空间$L$，$B$参数化余维$m+1$的中心$C$。考虑二步旗概形
\[
\mathcal F=\{(C,L)\in B\times G:C\subset L\},
\qquad p\colon\mathcal F\to G,\quad q\colon\mathcal F\to B.
\]
取$P=\mathbf P(E)$，$E=k^{N+1}$，以$A_G\subset E_G$和$A_B\subset E_B$分别表示定义$L$和$C$的秩$m$、$m+1$方程子丛。在可逆商的射影化约定下，旗的两种函子描述给出
\[
\mathcal F\simeq\mathbf P_G\bigl((E_G/A_G)^\vee\bigr),
\qquad
\mathcal F\simeq\mathbf P_B(A_B).
\]
前者选择$E_G/A_G$中的秩1局部直和因子，后者选择$A_B$的秩1商、以其核定义$L$。因此$p$确为相对维数$N-m$的射影丛，$q$确为相对维数$m$的射影丛；这不仅是逐点纤维的识别。特别地，$q$在$C$上的纤维为$Q(C)\simeq\mathbf P^m_{k(C)}$。故
\begin{align*}
\dim G&=m(N+1-m),\\
\dim B&=(m+1)(N-m),\\
\dim\mathcal F&=\dim G+N-m=\dim B+m.
\end{align*}
这些描述也可以逐个Grassmann仿射坐标图直接验证；Grassmann函子的商模构造及其仿射坐标图可对照\href{https://stacks.math.columbia.edu/tag/089R}{Stacks Project, Tag 089R}。

给闭集$Z\subset G$其约化闭子概形结构，设$\codim(Z,G)\ge q_0$。令$W=p^{-1}(Z)$。射影丛局部平凡，因此
\[
\dim W\le \dim\mathcal F-q_0=\dim B+m-q_0.
\]
把$W$分成有限个不可约分支。投影到$B$的非支配分支有闭的真像（因为$W\to B$固有），在$B$中除去这些像。对于每一个支配分支，其泛纤维维数为该分支维数减$\dim B$，至多$m-q_0$。有限型态射的泛纤维维数在一个非空开集上保持，或用固有态射纤维维数的上半连续性，可将$B$再缩到一个非空开集$U$，使所有几何纤维维数均至多$m-q_0$。参见\href{https://stacks.math.columbia.edu/tag/05F6}{Stacks Project, Tag 05F6}以及\href{https://stacks.math.columbia.edu/tag/0D4I}{Tag 0D4I}；其中关于纤维维数不受域扩张影响的结论保证了“几何纤维”的表述。

而$W_C=Q(C)\cap Z_{k(C)}$（交集在$k(C)$上理解），$\dim Q(C)=m$。故对$C\in U$有
\[
\codim\bigl(Q(C)\cap Z,Q(C)\bigr)\ge q_0.
\]
若$q_0>m$，支配分支不可能存在，缩小后交为空；约定空集余维为无穷即可统一表述。若基域有限，这个$U$的非空性不自动给出$U(k)$非空；引理的几何一般性和基域有理点存在性必须分开。
\end{analysisproof}

这个补证仅使用维数，不提供交集的不可约性。上一条的$Z=V(b^2-4ac)$已经是光滑、几何不可约的$G=\mathbf P^{2\vee}$中的闭子集，$Q(C)$为一般直线，交却有两个点。即使把原稿最后的“$Z$足够一般”解释为在平面二次曲线的参数空间中取一般光滑成员，一般直线仍与它相交于两个不同的点。因此正文不补造一个不可约性定理，而将可恢复的原句和不能成立的原因保存在V20-FAM012脚注。

\audit{V20-FAM-E04}{本段范围与未决事项}{第12节的文字回校结束；其未证推广没有被算作已证}
本轮末点为法稿物理73页末的手写“$Z$足够一般”句，英文48页末。法稿74页和英文49页均从第13节开始。第13节旧译逐字保留在新的继承模块中，本轮不声称已对其逐段回校。下一点为第13节标题及定义13.1。

第12节提出的任意$m$的8.7(c)型切丛判据、判别概形$Z$的全部性质、普通奇点推广及作者所引补篇的精确内部定位，仍须在后续统一审证。正文中字迹恢复、反例排错和本报告的余维补证各有明确范围；成功排版不能代替上述数学工作。

\subsection*{本段旧注全文存档：D17-M036--045}
\NewNote{以下为第十九版继承的旧注原文，用于追踪修订；其中“字迹不清”、旧对象归属和未决判断仅反映旧版状态。第二十版已辨明或订正的内容，以前面V20-FAM条目与当前正文为准。}

\paragraph*{旧注D17-M036}
这里按图像保留撇号$V'$、$V''$与上标$V^2$的差别。底本明确把$V^2$定义为分歧子预概形，不能仅凭字形把它认作首段的$V^{(2)}$。$X_{\operatorname{Grass}^{m}}$是底本此处的记号，不擅改成零点子概形$X^{(m)}$。关于光滑性及所引补篇的使用仍承接前文未完成的验证；“零点合并”保留为作者的解释性说法，而不是另立一个精确定理。

\paragraph*{旧注D17-M037}
底本在“it implies”后直接标“illegible”，故推论关系可能有缺文；这里保留缺口及作者“可以给出更精确陈述”的原话，不把它补写成已证明的唯一性或恢复映射。末句是原作者的写作交流，不是译者作出的后续承诺。

\paragraph*{旧注D17-M038}
“承认上述结果”承接前文仍待验证的推广讨论，不表示本译稿已经补齐证明；尤其不能据此恢复8.8等处的既有缺文。

\paragraph*{旧注D17-M039}
这里沿用本节$X^{(m)}$在格拉斯曼参数空间上的普遍线性截面族，$X_Q^{(m)}$表示将该族限制到$Q$。底本的“close subscheme”按语境译作“闭子概形”。

\paragraph*{旧注D17-M040}
底本给“regularly”加引号而未在本段重述精确定义；本译文保留这一限定，不把爆破同构或双有理性扩展为任意$C$均成立的结论。

\paragraph*{旧注D17-M041}
底本前两处只写$X^{(m)}$，但同段后文明确写$X_\xi^{(m)}$，并逐点使用剩余域，故正文补明纤维下标。底本大小写混用$K(\xi)$、$k(\xi)$，纤维处统一作$k(\xi)$。但全局判别集合$T$的几何不可约性究竟相对于哪个基域，不能从“$S=\operatorname{Spec}$”的残句和游移的$K(\xi)$唯一恢复，故不擅自补定。

\paragraph*{旧注D17-M042}
底本在“浸入”处写大写$F$，而本段随后仍以$f$表示同一态射，正文统一为$f$。底本又用$n$既表示$X$的维数，又作为$\phi_n$的指标；这里照录这一重用并指出，不自行断言两个用途必须取同一数值。作者后补的无分歧假设及预先替换$f$的要求仍是陈述的一部分，不能遗漏。

\paragraph*{旧注D17-M043}
底本此句交替使用$X^m$和$X^{(m)}$而未写参数下标；正文以“线性截面”表达其参数条件，原记号保留于此。本句仍是原作者提出的证明路径，不代表本译稿已经核实判别子概形的全部性质。

\paragraph*{旧注D17-M044}
底本在第一处交集中印作$Q(Z)\cap Z$，但本句定义的是由$C$决定的$Q(C)$，且紧接着在$Q(C)$中计算余维、后文也写$Q(C)\cap Z$，故正文将第一处$Q(Z)$订为$Q(C)$，原写法及依据保留于此。末句则确实写“if $Z$ is sufficiently general”，不擅自改成$C$。

\paragraph*{旧注D17-M045}
原句后半部为“also if $Z$ is geometrically irreducible [illegible, ask AG] [itou???] $Q(C)\cap Z$ if $Z$ is sufficiently general”。底本没有给出可辨的完整谓语；不凭“itou”可能表示“同样”就补出几何不可约性的结论，更不补写证明。
